\documentclass[10pt]{article}
\usepackage{helvet}
\renewcommand{\rmdefault}{phv} % Arial
\renewcommand{\sfdefault}{phv} % Arial
\usepackage{natbib}
\usepackage{bibentry}
\usepackage{amsmath,amsfonts,amstext,amssymb,amsthm}
\pagenumbering{gobble}

\setlength{\textwidth}{6.5in} 
\setlength{\oddsidemargin}{0in}  
\setlength{\evensidemargin}{0in} 
\setlength{\textheight}{8.5in}   
\setlength{\leftmargin}{1.1in} 
\setlength{\rightmargin}{1.0in} 
\setlength{\topmargin}{0in}      
\setlength{\headheight}{0in}     
\setlength{\headsep}{0in}        
\setlength{\footskip}{.5in}  


\usepackage{color}
\newcommand{\bl}{\textcolor{blue}}

\begin{document}

\bibliographystyle{apalike}
\nobibliography{lit_review_template}

\begin{center}
{\bf  {\large Literature review} \\
\vspace{12pt}
Clustering toward structured means with applications to high-dimensional genetic data}
\end{center}

\begin{enumerate}
 \item \bl{\bibentry{rojas2014change}}
    
    In this paper the authors prove some nice results about the fused lasso.  They consider the model
    \[
    Y_j = \mu_j + \varepsilon_j,\quad j = 1,\dots,p
    \]
    and study the fused lasso estimator of $\mu_1,\dots,\mu_p$ given by
    \[
    (\hat \mu_1,\dots,\hat \mu_p) = \underset{t_1,\dots,t_p}{\operatorname{argmin}} \sum_{j=1}^p (Y_j - t_j)^2 + \lambda \sum_{j=1}^{p-1}|t_{j+1} - t_j|.
    \]
    They do not consider any replication; that is, the sample is of size $1$.  They only observe a single observation $\mathbf{Y} \sim \text{Normal}(\boldsymbol{\mu},\sigma^2\cdot \mathbf{I}_p)$ with which they try to estimate $\boldsymbol{\mu} = (\mu_1,\dots,\mu_p)^T$.\\
    
    They also consider the multivariate case, when each $Y_j$ and each $\mu_j$ is a $d \times 1$ vector.\\
    
    The paper establishes conditions under which the fused lasso can be expected, for large $p$, to estimate change points with the correct sign within a certain distances of the true change points. They introduce the concept of $\varepsilon$-sign consistency: Basically, as $p \to \infty$, it is impossible for the fused lasso to recover the true change points exactly.  But this makes sense, because as the data are observed at a finer and finer resolution, the target of the true change point becomes smaller and smaller; that is, the probability of selecting exactly the true change point when a super fine grid of candidate change points is considered, vanishes.  However, these authors, with the concept of $\varepsilon$-sign consistency, consider the probability that a change point identified by the fused lasso will lie within some small distance of the true change point.  This is the thrust of the paper.
    
    
    \item \bl{\bibentry{guo2010pairwise}}
    
    The goal of this paper is quite similar to ours in that the authors wish to do clustering of high-dimensional profiles.  The authors assume the hierarchical model
    \begin{align*}
    \mathbf{Y}_i | \delta_{i1},\dots,\delta_{iK} &\sim f(\mathbf{y}_i|\delta_{i1},\dots,\delta_{iK}) = \prod_{k=1}^K \left[\phi(\mathbf{y}_i;\boldsymbol{\mu}_{k},\boldsymbol{\Sigma}_k)\right]^{\delta_{ik}}\\
    \delta_{i1},\dots,\delta_{iK} &\sim p(\delta_{i1},\dots,\delta_{iK}) = \pi^{\delta_{i1}} \cdots \pi^{\delta_{ik}} \cdot \mathbf{1}(\delta_{i1},\dots,\delta_{iK} \in  \{0,1\}, \sum_{k=1}^K \delta_{ik} = 1),
    \end{align*}
    giving
    \[
    \mathbf{Y}_i \sim \sum_{k=1}^K \pi_k \phi(\mathbf{y}_i;\boldsymbol{\mu}_{k},\boldsymbol{\Sigma}_k),
    \]
    where $\phi(\cdot ;\boldsymbol{\mu}_{k},\boldsymbol{\Sigma}_k)$ is the density of the multivariate Gaussian distribution with mean vector $\boldsymbol{\mu}_{k}$ and covariance matrix $\boldsymbol{\Sigma}_k$.  So the authors take a mixture distribution approach.  They then use maximum likelihood to estimate the parameters, but with a penalty which encourages some of the $K$ mean vectors to be equal to each other.   They propose to estimate the unknown parameters by maximizing the function
    \[
    \sum_{i=1}^n \log\left\{\sum_{k=1}^K \pi_k \phi(\mathbf{y}_i;\boldsymbol{\mu}_{k},\boldsymbol{\Sigma}_k), \right\} - \lambda \sum_{j=1}^p \sum_{1 \leq k , k' \leq K}|\mu_{k,j} - \mu_{k',j}|.
    \]
    The penalty term penalizes pairwise differences between the entries of the mean vectors $\boldsymbol{\mu}_1,\dots,\boldsymbol{\mu}_K$.  The authors use an EM-algorithm to maximize the objective function; they apply the penalty to the M-step of the ordinary EM-algorithm for mixture models.\\
    
    There is no penalization that encourages the means to have any particular structure.  The penalization in this paper really just controls the number of clusters.   So a large value of $K$ could be chosen and the $\lambda$  could be adjusted to choose the number of clusters.  Along the solution path (collection of maximizers of the objective function over different values of the tuning parameter $\lambda$), different numbers of clusters will be found---fewer clusters for larger values of $\lambda$.\\
    
    Our approach is very different in that we intend to fix the number of clusters and penalize the cluster means towards some structure (having few change points or being smooth).

    \end{enumerate}

\end{document}